\begin{problem}{과민성 대장 증후군}{standard input}{standard output}{1 second}{256 megabytes}

상원이는 과민성 대장 증후군을 앓고 있다. 과민성 대장 증후군의 원인은 스트레스! 상원이의 $N$일 동안의 스트레스 변화량 $A_1, \cdots, A_N$이 주어진다. $A_i$가 양수이면 $i$번째 날에 $A_i$만큼 스트레스가 쌓이고, 음수이면 $i$번째 날에 $-A_i$만큼 스트레스가 해소된다. 단, 변화를 관찰하기 시작한 시점의 스트레스 양은 $0$이며, 누적된 스트레스 양보다 해소하는 스트레스 양이 더 많을 경우 스트레스는 $0$이 될 때까지만 감소한다.

상원이는 스트레스가 $M$ 이상 쌓인 날에 복통을 겪게 될 때, 상원이가 며칠 동안 복통에 시달리게 되는지 알아보자.

\InputFile
첫째 줄에 스트레스 변화를 관찰한 일수 $N (1 \leq N \leq 100,000)$과 복통이 시작되는 스트레스의 양 $M (1 \leq M \leq 1,000,000,000)$이 주어진다.

둘째 줄에 스트레스 변화량 $a_i (-10,000 \leq a_i \leq 10,000)$가 공백으로 구분되어 $N$개 주어진다.

\OutputFile
상원이가 복통을 겪게 되는 일수를 출력한다.

\Examples

\begin{example}
\exmpfile{example.01}{example.01.a}%
\exmpfile{example.02}{example.02.a}%
\end{example}

\end{problem}

